% Part: second-order-logic
% Chapter: syntax-and-semantics
% Section: satisfaction

\documentclass[../../../include/open-logic-section]{subfiles}

\begin{document}

\olfileid[pt]{sol}{syn}{sat}

\olsection{Satisfação}

\begin{explain}
Para definir a relação de satisfação $\Sat{M}{!A}[s]$ para !!{formula}s de
segunda ordem, precisamos estender as definições de modo a abranger
!!{variable}s de segunda ordem. A noção de !!a{structure} é a mesma para a
lógica de segunda ordem e para a lógica de primeira ordem. A única diferença
está nas atribuições de variáveis~$s$: agora elas devem fornecer valores não
apenas para as !!{variable}s de primeira ordem, mas também para as
!!{variable}s de segunda ordem.
\end{explain}

\begin{defn}[Atribuição de variáveis]
Uma \emph{atribuição de variáveis}~$s$ para !!a{structure}~$\Struct{M}$ é
uma função que associa cada
\begin{enumerate}
\item !!{variable} de objeto~$\Obj{v_i}$ a um elemento de~$\Domain M$,
  O.s.\@, $s(\Obj{v_i}) \in \Domain{M}$;
\item variável de relação de $n$ argumentos~$\Obj{V_i^n}$ a uma relação de
  $n$ argumentos em~$\Domain{M}$, O.s.\@,
  $s(\Obj{V_i^n}) \subseteq \Domain{M}^n$;
\item variável de função de $n$ argumentos~$\Obj{u_i^n}$ a uma função de
  $n$ argumentos de $\Domain{M}$ para $\Domain{M}$, O.s.\@,
  $s(\Obj{u_i^n})\colon \Domain{M}^n \to \Domain{M}$;
\end{enumerate}
\end{defn}

\begin{explain}
!!^a{structure} atribui !!a{value} a cada !!{constant} e !!{function}, e
uma atribuição de variáveis de segunda ordem associa objetos e funções a
cada variável de objeto e de função. Juntas, elas nos permitem atribuir um
valor a cada termo.
\end{explain}

\begin{defn}[\usetoken{S}{value} de um termo]
Se $t$ é um termo da linguagem~$\Lang L$, $\Struct M$ é !!a{structure}
para~$\Lang L$, e $s$ é !!a{variable} atribuição para~$\Struct M$, o
\emph{!!{value}}~$\Value{t}{M}[s]$ é definido como para os termos de
primeira ordem, com a seguinte cláusula adicional:
\begin{quote}
\indcase{t}{\Atom{u}{t_1, \ldots, t_n}}{
\[
\Value{\indfrm}{M}[s] = s(u)(\Value{t_1}{M}[s], \ldots,
\Value{t_n}{M}[s]).
\]}
\end{quote}
\end{defn}

\begin{defn}[Variante em $x$]
Se $s$ é !!a{variable} atribuição para !!a{structure}~$\Struct M$, então
qualquer atribuição de !!{variable} $s'$ para $\Struct M$ que difira de $s$
no máximo no que atribui a $x$ é chamada \emph{variante em $x$} de~$s$. Se
$s'$ é uma variante em $x$ de $s$, escrevemos $\varAssign{s'}{s}{x}$.
(Analogamente para variáveis de segunda ordem $X$ ou $u$.)
\end{defn}

\begin{defn}
  Se $s$ é !!a{variable} atribuição para !!a{structure}~$\Struct M$ e
  $m \in \Domain{M}$, então a atribuição~$\Subst{s}{m}{x}$ é a atribuição
  de !!{variable} definida por
  \[\Subst{s}{m}{y} = \begin{cases}
    m & \text{se } y \ident x\\
    s(y) & \text{caso contrário},
  \end{cases}\]
  Se $X$ é uma !!{variable} de relação de $n$ argumentos e $M \subseteq
  \Domain{M}^n$, então $\Subst{s}{M}{X}$ é a atribuição de !!{variable}
  definida por
  \[\Subst{s}{M}{y} = \begin{cases}
    M & \text{se } y \ident X\\
    s(y) & \text{caso contrário}. 
  \end{cases}\]
  Se $u$ é uma !!{variable} de função de $n$ argumentos e
  $f\colon \Domain{M}^n \to \Domain{M}$, então $\Subst{s}{f}{u}$ é a
  atribuição de !!{variable} definida por
  \[\Subst{s}{f}{y} = \begin{cases}
    f & \text{se } y \ident u\\
    s(y) & \text{caso contrário}.
  \end{cases}\]
  Em cada caso, $y$ pode ser qualquer !!{variable} de primeira ou de
  segunda ordem.
\end{defn}

\begin{defn}[Satisfação]
Para !!{formula}s de segunda ordem~$!A$, a definição de satisfação é como
em \olref[fol][syn][sat]{defn:satisfaction}, com o acréscimo de:
\begin{enumerate}
\item \indcase{!A}{\Atom{X^n}{t_1, \dots, t_n}}{$\Sat{M}{\indfrm}[s]$
  sse $\langle \Value{t_1}{M}[s], \dots, \Value{t_n}{M}[s] \rangle \in
  s(X^n)$.}
\tagitem{prvAll}{%
  \indcase{!A}{\lforall[X][!B]}{$\Sat{M}{\indfrm}[s]$ sse, para todo $M
  \subseteq \Domain{M}^n$, $\Sat{M}{!B}[\Subst{s}{M}{X}]$.}}{}

\tagitem{prvEx}{%
  \indcase{!A}{\lexists[X][!B]}{$\Sat{M}{\indfrm}[s]$ sse existe pelo menos
  um $M \subseteq \Domain{M}^n$ tal que
  $\Sat{M}{!B}[\Subst{s}{M}{X}]$.}}{}

\tagitem{prvAll}{%
  \indcase{!A}{\lforall[u][!B]}{$\Sat{M}{\indfrm}[s]$ sse, para toda
  $f\colon \Domain{M}^n \to \Domain{M}$,
  $\Sat{M}{!B}[\Subst{s}{f}{u}]$.}}{}

\tagitem{prvEx}{%
  \indcase{!A}{\lexists[u][!B]}{$\Sat{M}{\indfrm}[s]$ sse existe pelo menos
  uma $f\colon \Domain{M}^n \to \Domain{M}$ tal que
  $\Sat{M}{!B}[\Subst{s}{f}{u}]$.}}{}
\end{enumerate}
\end{defn}

\begin{ex}
  Considere a !!{formula} $\lforall[z][(\Atom{X}{z} \liff \lnot
    \Atom{Y}{z})]$. Ela não contém quantificadores de segunda ordem, mas
  contém as !!{variable}s de segunda ordem $X$ e~$Y$ (aqui entendidas como
  unárias). A !!{sentence} de primeira ordem correspondente
  $\lforall[z][(\Atom{P}{z} \liff \lnot \Atom{R}{z})]$ diz que tudo o que
  pertence à interpretação de~$P$ não pertence à interpretação de~$R$, e
  vice-versa. Em !!a{structure}, a interpretação de !!a{predicate}~$P$ é
  dada pela interpretação~$\Assign{P}{M}$. Mas, para !!{variable}s de
  segunda ordem como $X$ e~$Y$, a interpretação não é fornecida pela própria
  !!{structure}, e sim por !!a{variable} atribuição. Como a !!{formula} de
  segunda ordem não é !!a{sentence} (ela contém as !!{variable}s livres $X$
  e~$Y$), só é satisfeita relativamente a !!a{structure}~$\Struct{M}$ junto
  com !!a{variable} atribuição~$s$.

  $\Sat{M}{\lforall[z][(\Atom{X}{z} \liff \lnot \Atom{Y}{z})]}[s]$
  sempre que os !!{element}s de~$s(X)$ não são !!{element}s de~$s(Y)$, e
  vice-versa, O.s.\@, sse $s(Y) = \Domain{M} \setminus s(X)$. Por exemplo,
  tome $\Domain{M} = \{1, 2, 3\}$. Como não há !!{predicate}s,
  !!{function}s nem !!{constant}s envolvidos, apenas o domínio de
  ~$\Struct{M}$ é relevante. Para $s_1(X) = \{1, 2\}$ e
  $s_1(Y) = \{3\}$, temos $\Sat{M}{\lforall[z][(\Atom{X}{z}
      \liff \lnot \Atom{Y}{z})]}[s_1]$.

  Em contraste, se $s_2(X) = \{1, 2\}$ e $s_2(Y) = \{2, 3\}$,
  $\Sat/{M}{\lforall[z][(\Atom{X}{z} \liff \lnot
      \Atom{Y}{z})]}[s_2]$. Isso ocorre porque
  $\Sat{M}{\Atom{X}{z}}[\Subst{s_2}{2}{z}]$ (pois $2 \in \Subst{s_2}{2}{z}(X)$), mas
  $\Sat/{M}{\lnot \Atom{Y}{z}}[\Subst{s_2}{2}{z}]$ (pois também $2 \in
  \Subst{s_2}{2}{z}(Y)$).
\end{ex}

\begin{ex}
  $\Sat{M}{\lexists[Y][(\lexists[y][\Atom{Y}{y}] \land
  \lforall[z][(\Atom{X}{z} \liff \lnot \Atom{Y}{z})])]}[s]$ se existe
  um $N \subseteq \Domain{M}$ tal que
  $\Sat{M}{(\lexists[y][\Atom{Y}{y}] \land \lforall[z][(\Atom{X}{z}
  \liff \lnot \Atom{Y}{z})])}[\Subst{s}{N}{Y}]$. E esse é o caso para
  qualquer $N \neq \emptyset$ (de modo que
  $\Sat{M}{\lexists[y][\Atom{Y}{y}]}[\Subst{s}{N}{Y}]$) e, como no exemplo
  anterior, $N = \Domain{M} \setminus s(X)$. Em outras palavras,
  $\Sat{M}{\lexists[Y][(\lexists[y][\Atom{Y}{y}] \land
  \lforall[z][(\Atom{X}{z} \liff \lnot \Atom{Y}{z})])]}[s]$ sse
  $\Domain{M} \setminus s(X)$ não é vazio, O.s.\@, $s(X) \neq
  \Domain{M}$. Assim, a !!{formula} é satisfeita, p.ex., se
  $\Domain{M} = \{1, 2, 3\}$ e $s(X) = \{1, 2\}$, mas não se
  $s(X) = \{1, 2, 3\} = \Domain{M}$.

  Como a !!{formula} não é satisfeita quando $s(X) = \Domain{M}$, a
  !!{sentence}
  \[
  \lforall[X][\lexists[Y][(\lexists[y][\Atom{Y}{y}] \land
      \lforall[z][(\Atom{X}{z} \liff \lnot \Atom{Y}{z})])]]
  \]
  nunca é satisfeita: para qualquer !!{structure}~$\Struct{M}$, a
  atribuição~$s(X) = \Domain{M}$ tornará falsa a !!{sentence}. Por outro
  lado, a sentença
  \[
  \lexists[X][\lexists[Y][(\lexists[y][\Atom{Y}{y}] \land
      \lforall[z][(\Atom{X}{z} \liff \lnot \Atom{Y}{z})])]]
  \]
  é satisfeita relativamente a qualquer atribuição~$s$, pois sempre podemos
  encontrar $M \subseteq \Domain{M}$ tal que $M \neq \Domain{M}$ (p.ex., $M = \emptyset$).
\end{ex}

\begin{ex}
  A !!{sentence} de segunda ordem~$\lforall[X][\lforall[y][X(y)]]$ diz que
  toda relação de $1$ argumento, O.s.\@, toda propriedade, vale para todo objeto. Isso
  claramente nunca é verdadeiro, pois, em cada~$\Struct{M}$, para uma
  atribuição de variáveis~$s$ com $s(X) = \emptyset$ e
  $s(y) = a \in \Domain{M}$, temos $\Sat/{M}{X(y)}[s]$. Isso significa que
  $!A \lif \lforall[X][\lforall[y][X(y)]]$ é equivalente, na lógica de
  segunda ordem, a $\lnot !A$, isto é: $\Sat{M}{!A \lif
  \lforall[X][\lforall[y][X(y)]]}$ sse $\Sat{M}{\lnot !A}$. Em outras
  palavras, na lógica de segunda ordem podemos definir $\lnot$ usando
  $\lforall$ e~$\lif$.
\end{ex}

\begin{prob}
  Mostre que, na lógica de segunda ordem, $\lforall$ e~$\lif$ podem definir
  os outros conectivos:
  \begin{enumerate}
    \item Prove que, na lógica de segunda ordem, $!A \land !B$ é equivalente
    a $\lforall[X][(!A \lif (!B \lif \lforall[x][X(x)])\lif
    \lforall[x][X(x)])]$.
    \item Encontre uma fórmula de segunda ordem que use apenas $\lforall$ e
    $\lif$ e seja equivalente a $!A \lor !B$.
  \end{enumerate}
\end{prob}

\end{document}
